%! Simplifying Rational Expressions — Unit 5, Lesson 5.1 — Slides (Beamer)
\documentclass[aspectratio=169, 11pt]{beamer}
\usepackage{algebra2-beamer}   % provides \forestheader{} and \sectionlabel[color]{}
% Standard coordinate grid + axes: {xmin}{xmax}{ymin}{ymax}
\newcommand{\lgrid}[4]{%
  \draw[step=1, linegray!45, thin] (#1,#3) grid (#2,#4);
  \draw[->, thick, charcoal] (#1,0)--(#2,0) node[below right, font=\tiny]{$x$};
  \draw[->, thick, charcoal] (0,#3)--(0,#4) node[left, font=\tiny]{$y$};}

\begin{document}

% TITLE SLIDE (CourseName is not defined in beamer, so the name is literal)
{
\setbeamercolor{background canvas}{bg=forest}
\begin{frame}[plain]
  \vspace{2.8cm}\hspace{0.55cm}%
  \begin{minipage}{0.9\textwidth}
    {\color{goldacc}\bfseries\small LESSON 5.1}\\[0.5cm]
    {\color{white}\bfseries\LARGE Simplifying Rational Expressions \& Domain Restrictions}\\[1.2cm]
    {\color{forestmid}\small Algebra 2: Shepherd~$\cdot$~Unit 5: Rational Functions}
  \end{minipage}
\end{frame}
}

% HOOK
\begin{frame}[t, plain]
  \forestheader{Two expressions. One disagreement.}
  \sectionlabel{Hook}
  \begin{columns}[T]
    \begin{column}{0.44\textwidth}
      \[ A(x)=\frac{x^2-x-6}{x^2-9} \qquad B(x)=\frac{x+2}{x+3} \]
      \vspace{-6pt}
      \begin{itemize}\setlength{\itemsep}{4pt}
        \item $B$ is $A$ factored and cancelled.
        \item They agree everywhere\dots{} except one input.
      \end{itemize}
    \end{column}
    \begin{column}{0.54\textwidth}
      \centering
      \renewcommand{\arraystretch}{1.4}
      \begin{tabular}{|l|c|c|c|c|}
      \hline
      $x$ & $0$ & $1$ & $4$ & $3$ \\ \hline
      $A(x)$ & $\frac23$ & $\frac34$ & $\frac67$ & \textbf{undef.} \\ \hline
      $B(x)$ & $\frac23$ & $\frac34$ & $\frac67$ & $\frac56$ \\ \hline
      \end{tabular}
    \end{column}
  \end{columns}
  \vspace{6pt}
  \begin{block}{Today}
    Cancelling changes how an expression \textbf{looks}. It never gives back an input that was already
    thrown away.
  \end{block}
\end{frame}

% WHY CANCELLING WORKS
\begin{frame}[t, plain]
  \forestheader{Cancelling is one rule, not a habit}
  \sectionlabel{Why it works}
  \begin{columns}[T]
    \begin{column}{0.52\textwidth}
      \[ \frac{a\cdot c}{b\cdot c}=\frac{a}{b}\cdot\frac{c}{c}=\frac{a}{b}\cdot 1=\frac{a}{b} \]
      \begin{itemize}\setlength{\itemsep}{4pt}
        \item You are dividing out a factor equal to \textbf{1}.
        \item It works \emph{only} if $c$ is a \textbf{factor}, and only if $c\neq0$.
      \end{itemize}
    \end{column}
    \begin{column}{0.46\textwidth}
      \begin{block}{Factors vs.\ terms}
        \textbf{Factors} are multiplied --- you may divide them out. \\[3pt]
        \textbf{Terms} are added --- you may never divide them out.
      \end{block}
      \vspace{3pt}
      {\footnotesize Claim: $\dfrac{x+3}{3}=x$. \; Test $x=6$: \; $\dfrac{9}{3}=3$, but $x=6$.
      \\[3pt] One number ends the argument.}
    \end{column}
  \end{columns}
\end{frame}

% THE FOUR STEPS
\begin{frame}[t, plain]
  \forestheader{The four steps --- and the order matters}
  \sectionlabel[goldacc]{Procedure}
  \begin{columns}[T]
    \begin{column}{0.50\textwidth}
      \begin{enumerate}\setlength{\itemsep}{5pt}
        \item \textbf{Factor} top and bottom completely (GCF first).
        \item \textbf{List the restrictions} --- from the \textbf{original} denominator.
        \item \textbf{Divide out} the shared factors.
        \item \textbf{Write} the simplest form \emph{with} its restrictions.
      \end{enumerate}
    \end{column}
    \begin{column}{0.48\textwidth}
      \begin{block}{The anchor}
        \[ \frac{x^2-x-6}{x^2-9}=\frac{(x-3)(x+2)}{(x-3)(x+3)}=\frac{x+2}{x+3} \]
        \centering $x\neq 3,\;-3$
      \end{block}
      \vspace{3pt}
      {\footnotesize Step 2 comes \emph{before} step 3 for a reason: after cancelling, $x=3$ is
      invisible --- but it is still excluded.}
    \end{column}
  \end{columns}
\end{frame}

% CANCELLING NEVER RESTORES AN INPUT
\begin{frame}[t, plain]
  \forestheader{Cancelling never restores an input}
  \sectionlabel{The big idea}
  \begin{columns}[T]
    \begin{column}{0.46\textwidth}
      \centering
      \begin{tikzpicture}[scale=0.36]
        \lgrid{-8}{6}{-4}{5}
        \draw[navy, dashed, thick] (-3,-4)--(-3,5);
        \draw[navy, dashed, thick] (-8,1)--(6,1);
        \begin{scope}
          \clip (-8,-4) rectangle (6,5);
          \draw[very thick, forest, domain=-8:-3.25, samples=100, smooth]
            plot (\x,{((\x)+2)/((\x)+3)});
          \draw[very thick, forest, domain=-2.8:6, samples=100, smooth]
            plot (\x,{((\x)+2)/((\x)+3)});
        \end{scope}
        \draw[forest, very thick, fill=white] (3,0.833) circle (5pt);
      \end{tikzpicture}
    \end{column}
    \begin{column}{0.54\textwidth}
      \begin{block}{Two restrictions, two jobs}
        $x=-3$: factor \textbf{stays} $\Rightarrow$ \textbf{vertical asymptote} \\[3pt]
        $x=3$: factor \textbf{cancels} $\Rightarrow$ \textbf{hole} at $\left(3,\tfrac56\right)$
      \end{block}
      \vspace{3pt}
      {\footnotesize This is Lesson 5.0's cancel test, run from the algebra side. The hole's height
      comes from the \emph{simplified} form --- the original has no value there at all.}
    \end{column}
  \end{columns}
\end{frame}

% OPPOSITE BINOMIALS
\begin{frame}[t, plain]
  \forestheader{Opposite binomials: the hidden factor of $-1$}
  \sectionlabel[goldacc]{$\frac{a-b}{b-a}=-1$}
  \begin{columns}[T]
    \begin{column}{0.50\textwidth}
      \[ \frac{5-x}{x-5}=\frac{-1(x-5)}{x-5}=-1 \]
      \begin{itemize}\setlength{\itemsep}{4pt}
        \item Same terms, \emph{every} sign flipped $\Rightarrow$ factor out $-1$.
        \item $x+3$ and $3+x$ are \emph{equal} --- only subtraction cares about order.
      \end{itemize}
    \end{column}
    \begin{column}{0.46\textwidth}
      \begin{block}{Worked}
        \[ \frac{x^2-16}{4-x}=\frac{(x-4)(x+4)}{-(x-4)}=-(x+4) \]
        \centering $x\neq4$
      \end{block}
      \vspace{3pt}
      {\footnotesize Dropping that $-1$ is the single most common lost point on this topic.}
    \end{column}
  \end{columns}
\end{frame}

% EQUIVALENCE
\begin{frame}[t, plain]
  \forestheader{What ``equivalent'' actually means}
  \sectionlabel{A2.EO.1d}
  \begin{center}
    Two rational expressions are \textbf{equivalent} when they give the same value at every input
    \textbf{both} are defined for.
  \end{center}
  \vspace{4pt}
  \begin{columns}[T]
    \begin{column}{0.48\textwidth}
      \begin{block}{Check by factoring}
        $\dfrac{x^2-4x}{x^2-16}=\dfrac{x(x-4)}{(x-4)(x+4)}=\dfrac{x}{x+4}$
      \end{block}
    \end{column}
    \begin{column}{0.48\textwidth}
      \begin{block}{Check by testing}
        At $x=1$: \; $\dfrac{-3}{-15}=\dfrac15$ \; and \; $\dfrac{1}{5}$. \\[3pt]
        One match proves nothing; one \emph{mismatch} disproves everything.
      \end{block}
    \end{column}
  \end{columns}
  \vspace{5pt}
  \begin{center}\footnotesize
    They still disagree at $x=4$ --- undefined on the left, $\tfrac14$ on the right. Different
    \emph{domains}, same \emph{values}: still equivalent.
  \end{center}
\end{frame}

% ERROR GALLERY
\begin{frame}[t, plain]
  \forestheader{The four ways to lose a point today}
  \sectionlabel{Watch out}
  \begin{columns}[T]
    \begin{column}{0.48\textwidth}
      \begin{block}{Cancelling terms}
        $\dfrac{x^2+9}{x+3}\neq x+3$ \\[3pt]
        {\footnotesize Test $x=1$: $\tfrac{10}{4}$ vs.\ $4$. And $x^2+9$ is \emph{prime}.}
      \end{block}
      \vspace{3pt}
      \begin{block}{Restrictions from the simplified form}
        The cancelled restriction disappears --- and the hole with it.
      \end{block}
    \end{column}
    \begin{column}{0.48\textwidth}
      \begin{block}{Losing the $-1$}
        $\dfrac{3-x}{x^2-9}=-\dfrac{1}{x+3}$, not $\dfrac{1}{x+3}$.
      \end{block}
      \vspace{3pt}
      \begin{block}{No restrictions at all}
        The answer is the expression \textbf{and} its restriction list. Half an answer is half credit.
      \end{block}
    \end{column}
  \end{columns}
\end{frame}

% ROADMAP
\begin{frame}[t, plain]
  \forestheader{Where this goes}
  \sectionlabel{Connections}
  \textbf{Today:} factor, restrict, divide out --- and know exactly which inputs your simplified form
  quietly dropped.\\[8pt]
  \begin{itemize}\setlength{\itemsep}{3pt}
    \item \textbf{5.2--5.3} Multiply, divide, add \& subtract --- every one of them starts by factoring.
    \item \textbf{5.5} Graphing from the equation: the cancel test decides hole vs.\ wall.
    \item \textbf{5.6} Solving rational equations --- an ``extraneous solution'' is nothing but one of
          today's excluded values sneaking back in.
  \end{itemize}
  \vspace{6pt}
  \begin{center}
    \textbf{Cancelling never restores an input.} Say it once more on the way out.
  \end{center}
\end{frame}

\end{document}
